\documentclass{article}
\usepackage[utf8]{inputenc}
\usepackage{amsmath,amssymb}
\usepackage[most]{tcolorbox}
\usepackage{geometry}
\geometry{margin=2.5cm}

\newcommand{\step}[1]{\par\noindent\hangindent=3.5em\hangafter=1%
  \makebox[3.5em][l]{\textbf{#1.}}}
\newcommand{\steptag}[1]{\hfill{\small\textsf{[#1]}}}

\newtcolorbox{proofbox}[1]{
  colback=gray!5, colframe=gray!60,
  fonttitle=\bfseries, title={#1},
  breakable, enhanced,
  left=4pt, right=4pt, top=4pt, bottom=4pt,
  before skip=10pt, after skip=10pt
}

\begin{document}

\begin{proofbox}{Parameterized inversion-derivative transport: there exist...}
\step{1} Parameterized inversion-derivative transport: there exist C\_der\textasciicircum{}0, C\_ch\textasciicircum{}0, C\_pol\textasciicircum{}0, C\_grp\textasciicircum{}0 >= 1 and kappa\_der\textasciicircum{}0, kappa\_ch\textasciicircum{}0, kappa\_pol\textasciicircum{}0 in (0, 1/2] such that, for every def-stage1-polar-witness-data tuple W with C\_der >= C\_der\textasciicircum{}0, C\_ch >= C\_ch\textasciicircum{}0, C\_pol >= C\_pol\textasciicircum{}0, C\_grp >= C\_grp\textasciicircum{}0, 0 < kappa\_der <= kappa\_der\textasciicircum{}0, 0 < kappa\_ch <= kappa\_ch\textasciicircum{}0, and 0 < kappa\_pol <= kappa\_pol\textasciicircum{}0, for every finite-dimensional exact-unit epsilon\_r-C*-algebra, every delta > 0, every s in \{+1, -1\}, and every 0 < r <= delta satisfying C\_ch*(epsilon\_r + delta) <= kappa\_ch, C\_pol*(epsilon\_r + delta) <= kappa\_pol, C\_grp*epsilon\_r < delta - C\_pol*(epsilon\_r*delta + delta\textasciicircum{}2), C\_der*(epsilon\_r + r) <= kappa\_der, and (1 + epsilon\_r)*(1 + C\_ch*(epsilon\_r + delta))*r + C\_grp*epsilon\_r < 2*delta, writing u\_delta for the unique first component of the inverse of Pi\_delta: calU x B\textasciicircum{}\{calH\}\_delta(J) -> S\_delta := Pi\_delta(calU x B\textasciicircum{}\{calH\}\_delta(J)), Pi\_delta(U, H) = U bold-dot H, and g\_\{sJ\}: B\_\{2delta\}\textasciicircum{}\{icalH\}(0) -> B\_\{2delta\}\textasciicircum{}\{calH\}(0) for the unique C\textasciicircum{}1 map such that, for every A in B\_\{2delta\}\textasciicircum{}\{icalH\}(0), f\_\{sJ\}(A + g\_\{sJ\}(A)) = 0, where f\_\{sJ\}(B) = (1/2)*(((J + B\textasciicircum{}dagger) bold-dot (sJ)\textasciicircum{}dagger) bold-dot (sJ bold-dot (J + B)) - J), define chi\_s(A) = sJ bold-dot (J + A + g\_\{sJ\}(A)) and the global C\textasciicircum{}1 map sigma(U) = u\_delta(U\textasciicircum{}dagger); then sigma maps chi\_s(B\_r\textasciicircum{}\{icalH\}(0)) into the same sJ-graph chart and, with F\_s(A) = phi\_\{sJ\}\textasciicircum{}par(sigma(chi\_s(A))), one has ||D(F\_s - id)(A) + 2*I\_\{icalH\}|| <= C\_der*(epsilon\_r + r) for every A in B\_r\textasciicircum{}\{icalH\}(0). \steptag{validated} \\[6pt]
\step{1.1} Joint universal thresholds: choose witnesses (C\_der\textasciicircum{}I,C\_ch\textasciicircum{}I,C\_pol\textasciicircum{}I,C\_grp\textasciicircum{}I,kappa\_der\textasciicircum{}I,kappa\_ch\textasciicircum{}I,kappa\_pol\textasciicircum{}I) from lem-stage1-inversion-derivative-control, (C\_pol\textasciicircum{}P,kappa\_pol\textasciicircum{}P) from lem-stage1-polar-retraction, (C\_ch\textasciicircum{}G,kappa\_ch\textasciicircum{}G) from lem-stage1-unitary-graph-control, and (C\_grp\textasciicircum{}A,C\_pol\textasciicircum{}A,kappa\_pol\textasciicircum{}A) from lem-stage1-approximate-group-laws. Define C\_der\textasciicircum{}0=C\_der\textasciicircum{}I, C\_ch\textasciicircum{}0=max(C\_ch\textasciicircum{}I,C\_ch\textasciicircum{}G), C\_pol\textasciicircum{}0=max(C\_pol\textasciicircum{}I,C\_pol\textasciicircum{}P,C\_pol\textasciicircum{}A), C\_grp\textasciicircum{}0=max(C\_grp\textasciicircum{}I,C\_grp\textasciicircum{}A), kappa\_der\textasciicircum{}0=kappa\_der\textasciicircum{}I, kappa\_ch\textasciicircum{}0=min(kappa\_ch\textasciicircum{}I,kappa\_ch\textasciicircum{}G), and kappa\_pol\textasciicircum{}0=min(kappa\_pol\textasciicircum{}I,kappa\_pol\textasciicircum{}P,kappa\_pol\textasciicircum{}A). These are universal, the four coefficient thresholds are at least 1, and the three margin thresholds lie in (0,1/2]. \steptag{validated} \\[4pt]
\step{1.2} Guard transport: for every W and every algebra, delta,s,r satisfying the hypotheses of node 1 with the thresholds of node 1.1, monotonicity transports the stated five guards simultaneously to the witnesses of lem-stage1-inversion-derivative-control and transports the relevant polar, graph, and group guards to lem-stage1-polar-retraction, lem-stage1-unitary-graph-control, and lem-stage1-approximate-group-laws. \steptag{validated} \\[6pt]
\step{1.2.1} Coefficient-margin guards: put e=epsilon\_r, x=e+delta, and z=e+r. By the definition of an epsilon\_r-C*-algebra e>=0, while delta>0 and r>0, so x,z>=0. Since C\_ch\textasciicircum{}I,C\_ch\textasciicircum{}G<=C\_ch\textasciicircum{}0<=C\_ch and kappa\_ch<=kappa\_ch\textasciicircum{}0<=kappa\_ch\textasciicircum{}I,kappa\_ch\textasciicircum{}G, the root graph guard implies C\_ch\textasciicircum{}I*x<=kappa\_ch\textasciicircum{}I and C\_ch\textasciicircum{}G*x<=kappa\_ch\textasciicircum{}G. The identical max/min argument gives C\_pol\textasciicircum{}q*x<=kappa\_pol\textasciicircum{}q for q in \{I,P,A\}, and C\_der\textasciicircum{}I*z<=kappa\_der\textasciicircum{}I. \steptag{validated} \\[4pt]
\step{1.2.2} Group-defect guards: put y=e*delta+delta\textasciicircum{}2>=0. For q=I and q=A, C\_grp\textasciicircum{}q<=C\_grp\textasciicircum{}0<=C\_grp and C\_pol\textasciicircum{}q<=C\_pol\textasciicircum{}0<=C\_pol. Hence C\_grp\textasciicircum{}q*e<=C\_grp*e<delta-C\_pol*y<=delta-C\_pol\textasciicircum{}q*y. These are respectively the group-defect hypotheses of lem-stage1-inversion-derivative-control and lem-stage1-approximate-group-laws. \steptag{validated} \\[6pt]
\step{1.2.2.1} Bind e := epsilon\_r for this node. Because an epsilon\_r-C*-algebra has epsilon\_r >= 0 and delta > 0, y := epsilon\_r*delta + delta\textasciicircum{}2 >= 0. For each q in \{I,A\}, the threshold construction gives C\_grp\textasciicircum{}q <= C\_grp\textasciicircum{}0 <= C\_grp and C\_pol\textasciicircum{}q <= C\_pol\textasciicircum{}0 <= C\_pol. Hence C\_grp\textasciicircum{}q*epsilon\_r <= C\_grp*epsilon\_r < delta - C\_pol*y <= delta - C\_pol\textasciicircum{}q*y, where the first inequality uses epsilon\_r >= 0 and the last uses y >= 0. Substituting y = epsilon\_r*delta + delta\textasciicircum{}2 yields C\_grp\textasciicircum{}q*epsilon\_r < delta - C\_pol\textasciicircum{}q*(epsilon\_r*delta + delta\textasciicircum{}2), exactly the group-defect guard required by the q=I inversion-derivative-control witness and the q=A approximate-group-laws witness. \steptag{validated} \\[4pt]
\step{1.2.3} Radius-chart guard: C\_ch\textasciicircum{}I<=C\_ch and C\_grp\textasciicircum{}I<=C\_grp, while e,x,r are nonnegative. Therefore (1+e)*(1+C\_ch\textasciicircum{}I*x)*r+C\_grp\textasciicircum{}I*e <= (1+e)*(1+C\_ch*x)*r+C\_grp*e < 2*delta, which is the remaining radius-chart hypothesis of lem-stage1-inversion-derivative-control. The unchanged assumptions s in \{+1,-1\} and 0<r<=delta provide its other typed hypotheses. \steptag{validated} \\[4pt]
\step{1.3} Polar binder identification: under the transported lem-stage1-polar-retraction guard, Pi\_delta(U,H)=U bold-dot H is a C\textasciicircum{}1 diffeomorphism from calU times B\_delta\textasciicircum{}\{calH\}(J) onto the same image S\_delta used in node 1. Hence its inverse exists uniquely, and the u\_delta explicitly bound in node 1 is exactly the first inverse component supplied by lem-stage1-polar-retraction and consequently the polar inverse component used by the imported group and inversion statements. \steptag{validated} \\[6pt]
\step{1.3.1} Existence on the identical domain: the P-witness guard C\_pol\textasciicircum{}P*(epsilon\_r+delta)<=kappa\_pol\textasciicircum{}P permits direct application of lem-stage1-polar-retraction. Its domain is calU times B\_delta\textasciicircum{}\{calH\}(J), its formula is the same Pi\_delta(U,H)=U bold-dot H, and its codomain is its image, exactly S\_delta:=Pi\_delta(calU times B\_delta\textasciicircum{}\{calH\}(J)) in node 1. Therefore it supplies a bijective C\textasciicircum{}1 diffeomorphism and a unique two-component inverse on precisely the root domain and image. \steptag{archived} \\[4pt]
\step{1.3.2} Uniqueness of the first component: a bijection has only one inverse map. Hence the first component called u\_delta by the explicit definite description in node 1 equals the first component (also called u\_delta) supplied by lem-stage1-polar-retraction. Whenever lem-stage1-approximate-group-laws, lem-stage1-smooth-polar-inverse, lem-stage1-smooth-unitary-operations, or lem-stage1-inversion-derivative-control uses the polar inverse for this same Pi\_delta, uniqueness identifies that component with the root u\_delta; no additional polar map is introduced. \steptag{archived} \\[4pt]
\step{1.4} Graph binder identification: for s in \{+1,-1\}, sJ lies in calU and hence in calUbar\_delta. Under the transported graph guard, lem-stage1-unitary-graph-control at V=sJ supplies exactly one C\textasciicircum{}1 map g\_\{sJ\}:B\_\{2delta\}\textasciicircum{}\{icalH\}(0) to B\_\{2delta\}\textasciicircum{}\{calH\}(0) satisfying the displayed f\_\{sJ\} equation. Thus it is exactly the g\_\{sJ\} bound in node 1, and the resulting graph parametrization is exactly chi\_s(A)=sJ bold-dot (J+A+g\_\{sJ\}(A)). \steptag{validated} \\[6pt]
\step{1.4.1} The chart base is admissible: because s is the real scalar +1 or -1, conjugate linearity of dagger and J\textasciicircum{}dagger=J give (sJ)\textasciicircum{}dagger=sJ. Bilinearity of bold-dot and the exact-unit identity J bold-dot J=J give (sJ)\textasciicircum{}dagger bold-dot (sJ)=s\textasciicircum{}2 J=J and (sJ) bold-dot (sJ)=J, so sJ has the right inverse sJ and zero unitary defect. By def-approximate-unitary-space, sJ belongs to calU=calUbar\_0, hence to calUbar\_delta because delta>0. \steptag{archived} \\[4pt]
\step{1.4.2} Apply graph control and identify by uniqueness: the G-witness guard C\_ch\textasciicircum{}G*(epsilon\_r+delta)<=kappa\_ch\textasciicircum{}G and node 1.4.1 allow lem-stage1-unitary-graph-control with V=sJ. Its displayed f\_V becomes verbatim the f\_\{sJ\} of node 1, and it supplies a unique C\textasciicircum{}1 g\_\{sJ\} with the required domain, codomain, and equation. The root definite description must therefore select this map. Substitution into the graph point gives verbatim chi\_s(A)=sJ bold-dot (J+A+g\_\{sJ\}(A)). \steptag{archived} \\[4pt]
\step{1.5} Global regularity and map identification: the conclusions of lem-stage1-unitary-graph-control and lem-stage1-polar-retraction satisfy the antecedents of lem-stage1-smooth-unitary-atlas and lem-stage1-smooth-polar-inverse; together with the transported lem-stage1-approximate-group-laws conclusion they satisfy the antecedents of lem-stage1-smooth-unitary-operations. Therefore, for the same explicitly identified u\_delta, sigma(U)=u\_delta(U\textasciicircum{}dagger) is a globally defined smooth, hence C\textasciicircum{}1, self-map of calU, and the upgrades change neither its points nor its first derivative. \steptag{validated} \\[6pt]
\step{1.5.1} Smooth-atlas antecedent: lem-stage1-unitary-graph-control says that its unique C\textasciicircum{}1 graph charts cover calU and that at every graph point ||D\_\{A\textasciicircum{}perp\}f\_V-I\_calH||<1. The elementary Neumann-series criterion implies that each such D\_\{A\textasciicircum{}perp\}f\_V is invertible. Thus the two hypotheses of lem-stage1-smooth-unitary-atlas hold, so the same graph functions and charts are smooth and define a smooth embedded manifold, without changing any point or first derivative. \steptag{archived} \\[4pt]
\step{1.5.2} Smooth-polar antecedent: lem-stage1-polar-retraction supplies Pi\_delta on the stated domain as a bijective C\textasciicircum{}1 diffeomorphism onto the open set S\_delta, hence in particular as a bijective C\textasciicircum{}1 local diffeomorphism. Combining this with the smooth embedded atlas from node 1.5.1 satisfies lem-stage1-smooth-polar-inverse. Consequently the same Pi\_delta and the same set-theoretic inverse (u\_delta,h\_delta) identified in node 1.3 are smooth, with no point or first derivative changed. \steptag{archived} \\[4pt]
\step{1.5.3} Smooth operation: the A-witness polar and group-defect guards permit lem-stage1-approximate-group-laws, which defines sigma(U)=u\_delta(U\textasciicircum{}dagger) on all of calU for the same polar inverse identified in node 1.3. Its conclusion together with nodes 1.5.1 and 1.5.2 supplies all three antecedents named in lem-stage1-smooth-unitary-operations. That lemma upgrades this same sigma to a smooth map calU to calU and explicitly changes no point or first derivative. Therefore it is globally defined and C\textasciicircum{}1 exactly as required in node 1. \steptag{archived} \\[4pt]
\step{1.5.4} Atlas and polar upgrades: lem-stage1-unitary-graph-control gives a covering by the unique C\textasciicircum{}1 graphs and ||D\_\{A\textasciicircum{}perp\}f\_V-I\_calH||<1 at every graph point; the Neumann-series criterion makes each displayed derivative invertible, so lem-stage1-smooth-unitary-atlas applies without changing points or first derivatives. Lem-stage1-polar-retraction gives Pi\_delta as a bijective C\textasciicircum{}1 diffeomorphism onto open S\_delta, hence a bijective C\textasciicircum{}1 local diffeomorphism; with that smooth atlas, lem-stage1-smooth-polar-inverse applies and makes the same Pi\_delta and same inverse (u\_delta,h\_delta) smooth, again without changing points or first derivatives. \steptag{validated} \\[4pt]
\step{1.5.5} Operation upgrade: the transported polar and group-defect guards apply lem-stage1-approximate-group-laws, defining sigma(U)=u\_delta(U\textasciicircum{}dagger) on all calU for the unique inverse component identified in node 1.3. This result and node 1.5.4 are precisely the three antecedents named by lem-stage1-smooth-unitary-operations: approximate group laws, smooth unitary atlas, and smooth polar inverse. Therefore that external makes the same sigma smooth as a self-map of calU and explicitly changes no point or first derivative; in particular sigma is the global C\textasciicircum{}1 map required by node 1. \steptag{validated} \\[4pt]
\step{1.6} Final derivative transport: apply lem-stage1-inversion-derivative-control with its selected I-witnesses and the transported five guards. After the polar, graph, and sigma identifications above, its chi\_s, sigma, F\_s and chart are precisely those in node 1. It gives same-chart retention and norm at most C\_der\textasciicircum{}I*(epsilon\_r+r), which is at most C\_der*(epsilon\_r+r) because C\_der>=C\_der\textasciicircum{}0=C\_der\textasciicircum{}I and epsilon\_r+r is nonnegative. This is exactly the conclusion of node 1. \steptag{validated} \\[6pt]
\step{1.6.1} Direct producer application: the I-witness graph, polar, group, derivative, and radius-chart guards established in node 1.2 meet every hypothesis of lem-stage1-inversion-derivative-control for the fixed algebra, delta,s,r. By nodes 1.3-1.5, every anaphoric polar inverse, graph function, chi\_s, and global sigma in that external is the explicitly bound map of node 1. Therefore the external conclusion says that this sigma maps chi\_s(B\_r\textasciicircum{}\{icalH\}(0)) into the same sJ-graph chart and that the resulting F\_s(A)=phi\_\{sJ\}\textasciicircum{}par(sigma(chi\_s(A))) obeys ||D(F\_s-id)(A)+2*I\_\{icalH\}||<=C\_der\textasciicircum{}I*(epsilon\_r+r) for every A in that ball. \steptag{archived} \\[4pt]
\step{1.6.2} Enlarge only the displayed bound: epsilon\_r+r>=0 and the witness hypotheses give C\_der>=C\_der\textasciicircum{}0=C\_der\textasciicircum{}I. Multiplication by epsilon\_r+r preserves order, so C\_der\textasciicircum{}I*(epsilon\_r+r)<=C\_der*(epsilon\_r+r). Combining this with node 1.6.1 preserves the exact chart-retention conclusion and yields the precise inequality asserted in node 1. \steptag{archived} \\[4pt]
\end{proofbox}

\end{document}
